\documentclass[10pt,a4paper]{amsart}
\usepackage[margin=19mm]{geometry}
\usepackage{amsmath,amssymb,mathtools}
\usepackage[scr=rsfs,cal=euler]{mathalfa}
\usepackage{xfrac}
\usepackage[unicode,hidelinks,bookmarks=false]{hyperref}
\newcommand{\ie}{i.e.\ }
\newcommand{\cf}{cf.\ }
\newcommand{\resp}{respectively\ }
\newcommand{\Subsection}{\subsection}
\providecommand{\nobreakdash}{\nobreak\hbox{-}}
\setlength{\emergencystretch}{2em}
\multlinegap0pt
\usepackage{bm}
\usepackage{bbm}
\usepackage{dsfont}
\usepackage{xspace}
\usepackage{cancel}
\usepackage{mathtools}
\usepackage{amsthm}
\makeatletter
\@ifundefined{lemma}{\newtheorem{lemma}{Lemma}}{}
\makeatother
\title[Integer moments of the derivatives of the Riemann zeta funct]{Integer moments of the derivatives of the Riemann zeta function}
\author{Christopher Hughes, Andrew Pearce-Crump}
\date{Source: arXiv:2509.07792v1 (CC BY 4.0)}
\begin{document}
\maketitle

\noindent\emph{Source and licence.} Integer moments of the derivatives of the Riemann zeta function. Version arXiv:2509.07792v1 (CC BY 4.0), distributed under the Creative Commons Attribution 4.0 International licence, as recorded in the arXiv licence metadata for this exact version and shown on the abstract page https://arxiv.org/abs/2509.07792v1. Full rights notice: licenses/arxiv-2509.07792-CC-BY-4.0.txt.\par\smallskip

\noindent\textit{Editorial scope.} Counted unit: the Lemma whose source label is \texttt{lem:vandermonde} in the article (source0.tex lines 360--372, source label \texttt{lem:vandermonde}), delivered together with the article's own complete original proof of it (source0.tex lines 374--447, one \texttt{proof} environment of 74 lines). No mathematics is written, repaired or re-derived: statement and proof are the article's own text line for line. Required context retained uncounted: none. Every definition, display and label of those lines is retained unchanged. Omitted from this package and not redistributed: 1--359, 373--373, 448--1218. Nothing inside a retained excerpt is omitted or altered: every retained body is delivered exactly as the article's lines are written, so no omission receipt of any kind accompanies this package. Citations: the retained excerpts carry no citation command at all, so no bibliography entry of the article is transcribed and no reference is redistributed. Reference adaptation: none was needed and none was made; every reference inside the delivered unit resolves inside it, so no unresolved reference or citation is produced. Printed slips kept literal: no printed word, symbol, hypothesis or number of the counted statement or of its proof was altered. Layout and class: the article's own class is \texttt{amsart} with packages including fouriernc, geometry, xcolor, graphicx, placeins, natbib; the wrapper is the lane's pinned \texttt{amsart} editorial wrapper at 10pt on A4, which restates the theorem-like environments the retained text needs and adds the article's own macro definitions that the retained text needs (none), with extra wrapper packages (bm, bbm, dsfont, xspace, cancel, mathtools). The delivered numbering is the unit's own. Assets: no retained span contains a figure environment, a graphics-inclusion command, a PDF-page inclusion, a table or a TikZ picture; no image file is copied into this package, the package declares no asset dependency and no image file is redistributed with it. Licence: version arXiv:2509.07792v1 of this article is distributed under the Creative Commons Attribution 4.0 International licence, as recorded in the arXiv licence metadata for this exact version and shown on the abstract page https://arxiv.org/abs/2509.07792v1. A full rights notice is delivered at licenses/arxiv-2509.07792-CC-BY-4.0.txt. Characters: every retained excerpt is pure ASCII, so the delivered engine, XeLaTeX with its default Latin Modern text font, reports no missing character.
\begin{lemma}\label{lem:vandermonde}
For $k\in \mathbb{N}$ and $n\in \mathbb{Z}$, if $\alpha_1,\dots,\alpha_k \neq 0$ then
\begin{multline*}
\sum_{\ell=1}^k (-1)^{k+\ell}   \alpha_\ell^{n-1}  \Delta_{k-1}\left(\alpha_1,\dots,\alpha_k ; \alpha_\ell\right)\\
=
\begin{cases}
\displaystyle \Delta_{k}\left(\alpha_1,\dots,\alpha_{k}\right) h_{n-k}\left(\alpha_1, \dots \alpha_k\right) & \text{ if } n \geq k \\
0 & \text{ if } 1\leq n \leq k-1 \\
\noalign{\vskip1ex}
\displaystyle (-1)^{k+1} \Delta_{k}\left(\alpha_1,\dots,\alpha_{k}\right)   \frac{ h_{|n|}\left(\alpha_1^{-1}, \dots \alpha_k^{-1}\right)}{\alpha_1 \alpha_2 \dots \alpha_k} & \text{ if } n \leq 0.
\end{cases}
\end{multline*}
\end{lemma}

\begin{proof}
Let
\begin{equation*}
X = \sum_{\ell=1}^{k} (-1)^{k+\ell} \alpha_\ell^{n-1} \Delta_{k-1}\left(\alpha_1,\dots,\alpha_k ; \alpha_\ell\right) .
\end{equation*}


Observe that $X$ can be expressed as a determinant,
\[
\renewcommand{\arraystretch}{1.4}
X = \det \begin{pmatrix}
1 & \alpha_1 & \alpha_1^2 & \dots & \alpha_1^{k-2} & \alpha_1^{n-1}  \\
1 & \alpha_2 & \alpha_2^2 & \dots & \alpha_2^{k-2} & \alpha_2^{n-1}  \\
\vdots & \vdots & \vdots & \dots & \vdots & \vdots\\
1 & \alpha_{k} & \alpha_{k}^2 & \dots & \alpha_{k}^{k-2}  & \alpha_k^{n-1}
\end{pmatrix}
\]
where the terms in the last column are replaced by $\alpha_j^{n-1}$ (in the standard Vandermonde these terms would be $\alpha_j^{k-1}$). Expanding the determinant down its last column, it is clear that it equals $X$.

Note that if $n\leq k-1$ then the last column will be equal to one of the other columns in the matrix, and thus the determinant will be zero.

To complete the proof of the lemma, we now evaluate $X$ in a  different manner, via row manipulations.

Subtract the first row from all subsequent rows. Pull out a factor of $(\alpha_j-\alpha_1)$ from the $j$\textsuperscript{th} row (for $j\geq 2$).

Then subtract the new second row from all subsequent rows, and pull out a factor of $(\alpha_j-\alpha_2)$ from the $j$\textsuperscript{th} row (for $j\geq 3$ this time).

Repeat this process all the way down to subtracting what remains in the the $k-1$\textsuperscript{st} row from what remains in the $k$\textsuperscript{th} row, and pulling out a factor of $(\alpha_k-\alpha_{k-1})$.

The terms that are factored out during this process are
\[
\prod_{\ell=1}^{k-1} \prod_{j=\ell+1}^k \left(\alpha_j-\alpha_\ell\right)  = \Delta_{k}\left(\alpha_1,\dots,\alpha_{k}\right)
\]
and the remaining matrix determinant is
\[
\det
\begin{pmatrix}
A_{1,1} & A_{1,2} & \dots & A_{1,k-1} & B_{1} \\
A_{2,1} & A_{2,2} & \dots & A_{2,k-1} & B_{2} \\
\vdots & \vdots & \vdots & \vdots & \vdots\\
A_{k,1} & A_{k,2} & \dots & A_{k,k-1} & B_{k} \\
\end{pmatrix}
\]
where
\[
A_{\ell,m} = \sum_{\substack{n_1,\dots,n_{\ell} \geq 0 \\ n_1 + \dots +n_\ell  = m-\ell}} \alpha_1^{n_1} \dots \alpha_{\ell}^{n_\ell}
\]
and where $B_\ell = A_{\ell,n}$ in the case $n \geq 1$, and in the case $n \leq 0$,
\[
B_{\ell} = (-1)^{\ell+1} \sum_{\substack{n_1,\dots,n_{\ell} \geq 0 \\ n_1 + \dots +n_\ell = |n|}} \alpha_1^{-n_1-1} \dots \alpha_{\ell}^{-n_\ell-1}.
\]

Crucially note that $A_{\ell,m}=0$ for $m<\ell$ and $A_{\ell,\ell}=1$. This makes the determinant trivial to evaluate - the matrix is upper triangular, and so the determinant will equal $B_k$.

That is, we have shown that
\[
X = \Delta_{k}\left(\alpha_1,\dots,\alpha_{k}\right) B_k
\]
and this is the right-hand of the statement of the lemma, since if $1 \leq n\leq k-1$ we have
\[
B_k = A_{k,n} = 0
\]
and if $n \geq k$ we have
\begin{align*}
B_k &= A_{k,n} = \sum_{\substack{n_1,\dots,n_{k} \geq 0 \\ n_1 + \dots +n_k  = n-k}} \alpha_1^{n_1} \dots \alpha_{k}^{n_k} \\
&= h_{n-k}\left(\alpha_1,\dots,\alpha_k\right)
\end{align*}
and if $n \leq 0$ we have
\begin{align*}
B_k &= (-1)^{k+1} \sum_{\substack{n_1,\dots,n_k \geq 0 \\ n_1 + \dots +n_k = |n|}} \alpha_1^{-n_1-1} \dots \alpha_{k}^{-n_k-1} \\
&= \frac{ h_{|n|}\left(\alpha_1^{-1}, \dots \alpha_k^{-1}\right)}{\alpha_1 \alpha_2 \dots \alpha_k}
\end{align*}
as required.
\end{proof}

\end{document}
